In this section, we discuss the Chow $t$-structure on $\SHRat$ and how it can be used to define an interesting lax symmetric monoidal functor
\[\nur : \Sp_{C_2,i2} \to \SHRat.\]
In the context of $\mathrm{DM}(k)$, the Chow $t$-structure was first studied by Bondarko \cite{Bondarko}.\footnote{The publication history of this definition is somewhat complicated. In the published version of \cite{Bondarko}, the Chow $t$-structure is only conjectured, not shown, to exist. However, in a later version of \cite{Bondarko} which appeared on the arXiv, Bondarko gave a construction of the Chow $t$-structure.}
In the context of $\SHC^{\at}$ it was studied by Pstragowski in \cite{Pstragowski}.
The Chow $t$-structure on $\SH(k)$ is the subject of forthcoming work of Bachmann--Kong--Wang--Xu, who we thank for carefully explaining their results to us. Among their results is a determination of the heart of the Chow $t$-structure in full generality.\footnote{In the time since the first version of this work was posted, the theorems of Bachmann--Kong--Wang--Xu appeared in \cite{Chow}}
%The Chow $t$-structure on $\SH(k)$ is the subject of work of Bachmann--Kong--Wang--Xu \cite{Chow}---one of their results is a determination of the heart of the Chow $t$-structure in full generality.

%The Chow $t$-structure on $\SH(k)$ is the subject of work of Bachmann--Kong--Wang--Xu \cite{Chow}---one of their results is a determination of the heart of the Chow $t$-structure in full generality.

%who we thank for carefully explaining their results to us, among which is a determination of the heart of the Chow $t$-structure in full generality.

In \Cref{sec:chow}, we define the Chow $t$-structure on $\SHRat$, determine its heart and give a formula for the $t$-structure homotopy objects.
In \Cref{sec:nu}, we use the Chow $t$-structure to define the functor
\[\nur : \Sp_{C_2} \to \SHRat.\]
We prove several basic properties of this functor, compare it with the functor $\Gamma_*$ defined in \Cref{sec:top}, and show that several Artin-Tate $\R$-motivic spectra may be recovered from their Betti realizations via $\nur$.
%
%The other filtration that we discuss is Voevodsky's effective slice filtration \cite{VoeOpenProbs}.
%In \Cref{sec:eff}, we will show that the effective slice filtration of an $\R$-motivic spectrum admits a simple description on the level of our filtered model for $\SHRat$.
%As a consequence, we demonstrate that the $\ta$-Bockstein spectral sequence is essentially equivalent to Isaksen's effective spectral sequence, which is associated to the Betti realization of the effective slice tower.
%
%\begin{dfn}
%  Using [ref] we can construct $t$-structures on $\SHgc$ simply by specifying which full subcategory we wish to to be the $\geq 0$ part. This allows us to define several $t$-structures fairly easily,
%  \begin{itemize}
%  \item The homotopy $t$-structure has connective part generated under colimits by
%    \[ \{ \R, \C, (\mathbb{G}_m)^n\ |\ n \in \Z \} \]
%  \item The Chow $t$-structure has connective part generated under colimits by
%    \[ \{ \R, \C, (\mathbb{P}^1)^n\ |\ n \in \Z \} \]
%  \item The vertical $t$-structure has connective part generated under colimits by
%    \[ \{ \R, \C, (\mathbb{P}^1)^n\ |\ n \in \Z \} \]
%  \end{itemize}
%\end{dfn}
%
\subsection{The Chow $t$-structure} \label{sec:chow}\ 

In this section, we define the Chow $t$-structure on Artin-Tate $\R$-motivic spectra, compute its heart, and describe the $t$-structure homotopy objects, which we denote $\pi^{\ch} _k$.

\begin{dfn}
  By \cite[Proposition 1.2.1.16]{HA}, we can construct a $t$-structure on $\SHRat$ with
	\begin{itemize}
	\item connective part $(\SHRat)^{\cn} _{\geq 0}$ generated under colimits and extensions by \[\{ \Ss_2^{n+k_1,n+k_2,n}\ \vert \ k_1 \geq 0 \text{ and } k_1 + k_2 \geq 0 \}.\]
	\item $(-1)$-coconnective part $(\SHRat)^{\cn} _{< 0}$ spanned by those $X \in \SHRat$ with $\pi_{n+k_1, n+k_2, n} ^{\R} X = 0$ for all $n,k_1,k_2 \in \Z$ satisfying $k_1 \geq 0$ and $k_1+k_2 \geq 0$.
	\end{itemize}
  We call this $t$-structure the \emph{Chow $t$-structure} on $\SHRat$.
%
%  Note that the Chow $t$-structure is compatible with the symmetric monoidal structure in the sense that the full subcategory of Chow connective objects contains the unit and it closed uner the tensor product.
%
%  Since the functors $X \mapsto \pi_{p,q,w} ^{\R} X$ are jointly conservative on $\SHRat$, it is easy to see that the Chow $t$-structure is right complete.
  %simply by specifying which full subcategory we wish to to be the $\geq 0$ part.
%
  %We define the Chow $t$-structure to have connective part generated under colimits by
  %\[ \{ \o_\R, \Spec(\C), (\mathbb{P}^1)^n\ \vert \ n \in \Z \}. \]
  %Note that the full subcategory of Chow connective objects contains the unit and it closed uner the tensor product.
%
  Moreover, we let $\tau_{\geq 0} ^{\cn} : \SHRat \to (\SHRat)^{\cn} _{\geq 0}$ denote the connective cover with respect to the Chow $t$-structure, and let $\pi_n ^{\ch} : \SHRat \to \SHR_{i2}^{\at, \ch}$ denote the homotopy object functors.
\end{dfn}

The following result is immediate from the definition:
\begin{prop}
  The Chow $t$-structure on $\SHRat$ is compatible with the symmetric monoidal structure, right complete and compatible with filtered colimits.
  Moreover, $\Sigma^{n,n,n}$ restricts to an equivalence $\Sigma^{n,n,n} : (\SHRat)^{\cn} _{\geq 0} \to (\SHRat)^{\cn} _{\geq 0}$.
  %Moreover, for any $X \in \SHRat$ there are natural equivalences
  %\begin{align*}
    %\tau_{\geq 0} ^{\cn} (\Sigma^{n,n,n} X) \simeq \Sigma^{n,n,n} \tau_{\geq 0} ^{\cn} (X) \\
    %\pi{0} ^{\ch} (\Sigma^{n,n,n} X) \simeq \Sigma^{n,n,n} \tau_{\geq 0} ^{\cn} (X) 
  %\end{align*}
\end{prop}

The main two theorems that we prove in this section are as follows:

\begin{thm} \label{thm:chow-heart}
  Let $\SHRatch$ denote the heart of $\SHRat$ with respect to the Chow $t$-structure. Then there is a symmetric monoidal equivalence of categories
  \[\SHRatch \simeq \Comod(\uAbh_{i2}; \underline{(\MU_2)_{2*} \MU_2}),\]
  where $\Comod(\uAbh; \underline{(\MU_2)_{2*} \MU_2})$ is the category of evenly-graded $(\underline{(\MU_2)_{*}}, \underline{(\MU_2)_{*} \MU_2})$-comodules in $C_2$-Mackey functors.
\end{thm}

\begin{thm} \label{thm:pi0-MGL}
  Given any $X \in \SHRat$ and $n,k \in \Z$, there is an isomorphism
  \[(\pi_k ^{\ch} X)_n \cong \underline{(\MGL_2)}_{n+k,n,n} (X).\]
\end{thm}

\begin{rec}
In \Cref{thm:Cta}, we constructed a symmetric monoidal equivalence
\[ \Mod(\SHRat; C\ta) \simeq \Mod(\Sp_{C_2, i2}; \uZt) \otimes_{\Z} \IndCoh(\Mfg). \]
Moreover, in \Cref{lem:Cta-t-structure} we equipped $\Mod(\Sp_{C_2, i2}; \uZt) \otimes_{\Z} \IndCoh(\Mfg)$ with a $t$-structure whose heart is identified with $\Comod(\uAbh; \underline{(\MU_2)_{2*} \MU_2})$.
  We call the induced $t$-structure on $\Mod(\SHRat; C\ta)$ the \emph{tensor $t$-structure} and write $\Mod(\SHRat; C\ta)^{\th}$ for the heart.
\end{rec}

To prove \Cref{thm:chow-heart}, it will suffice to prove the following:

\begin{prop} \label{prop:t-compare}
  The tensor $t$-structure on $\Mod(\SHRat; C\ta)$ is induced by the Chow $t$-structure on $\SHRat$: a $C\ta$-module $X$ is tensor (co)connective if and only if its underlying Artin-Tate $\R$-motivic spectrum is Chow (co)connective.

  Moreover, the induced symmetric monoidal functor
  \[\Mod(\SHRat; C\ta)^{\th} \to \SHR_{i2}^{\at, \ch}\]
  is an equivalence of categories.
  %Equip $\SHRat$ with the Chow $t$-structure and $\Mod(\SHRat; C\ta)$ with the tensor $t$-structure.
  %Then the forgetful functor
  %\[\Mod(\SHRat; C\ta) \to \SHRat\]
  %is $t$-exact and induces an equivalence upon taking the heart.
\end{prop}

%In fact, we will even show in \Cref{lem:tensor-t-ta} that the tensor $t$-structure on is induced by the Chow $t$-structure.

First, we need a lemma:

\begin{lem} \label{lem:Cta-pi0}
  We have $C\ta \in \SHRatch$ and the unit map $\Ss_2 ^{0,0,0} \to C\ta$ induces an equivalence $\pi_0 ^{\ch} \Ss_2 ^{0,0,0} \simeq C\ta$.
\end{lem}

\begin{proof}
  First, we note that $\Ss_2^{0,0,-1} \geq 1$ in the Chow $t$-structure.
  Since $$\Ss_2^{0,0,-1} \simeq \Sigma \Ss_2^{0,1,0} \otimes \Ss_2^{-1,-1,-1},$$ it suffices to show that $\Ss_2^{0,1,0} \geq 0$.
  This follows immediately from the cofiber sequence $\Sigma^{\infty} _+ \Spec(\C)_2 \to \Ss_2^{0,0,0} \to \Ss_2^{0,1,0}$.

  Combining this fact with the cofiber sequence $\Ss_2^{0,0,-1} \xrightarrow{\ta} \Ss_2^{0,0,0} \to C\ta$, we find that $C\ta \geq 0$ and that $\Ss_2^{0,0,0} \to C\ta$ induces an equivalence on $\pi_0 ^{\ch}$.

  To conclude, it suffices to show that $C\ta \leq 0$.
  By definition, we must show that $\pi_{n+k_1, n+k_2, n} ^{\R} C\ta = 0$ for all $n,k_1,k_2$ satisfying $k_1 > 0$ and $k_1 + k_2 > 0$.
  This follows directly from \Cref{thm:omnibus-vanishing}(2).
\end{proof}

We note down the following corollary for later use: 

\begin{cor} \label{cor:cta-homotopy}
  Suppose that $X \in \SHRat$ is a filtered colimit of Artin-Tate $\R$-motivic spectra that each admit a finite cell structure with all cells of the form $\Ss_2^{n,n,n}$.
  Then $X \otimes C\ta \in \SHRatch$ and $X \to X \otimes C\ta$ induces an equivalence $\pi_0 ^{\ch} X \to X \otimes C\ta$.
\end{cor}

\begin{proof}
  It is clear that the collection of $X$ which satisfy the conclusions of the corollary is closed under filtered colimits and extensions, so it suffices to assume that $X \simeq \Ss_2^{n,n,n}$
  Since $\Sigma^{n,n,n}$ is an automorphism of the Chow $t$-structure, we may reduce to the case of $X \simeq \Ss_2^{0,0,0}$, which is precisely \Cref{lem:Cta-pi0}.
\end{proof}

%\begin{lem} \label{lem:tensor-t-ta}
  %The tensor $t$-structure on $\Mod(\SHRat; C\ta)$ is induced by the Chow $t$-structure on $\SHRat$: a $C\ta$-module $X$ is tensor (co)connective if and only if it is Chow (co)connective.
%\end{lem}

\begin{proof}[Proof of \Cref{prop:t-compare}]
  We first show that the tensor $t$-structure is induced by the Chow $t$-structure.
  We begin by showing that the connective part of the tensor $t$-structure on $\Mod(\SHRat; C\ta)$ is generated under colimits and extensions by
  \[ \{ \Sigma^{n+k_1, n+k_2,n} C\ta \ \vert \ k_1 \geq 0 \text{ and } k_1 + k_2 \geq 0 \}. \]

  The $t$-structure on $\IndCoh(\Mfg)$ has connective part generated under colimits and extensions by $\omega^{\otimes n} _{\mathbb{G}/ \Mfg}$ for $n \in \Z$.
  On the other hand, the $t$-structure on $\Mod(\Sp_{C_2, i2}; \uZt)$ has connective part generated under colimits and extensions by $\Sigma^{k_1+ k_2 \sigma} \uZt$, where $k_1 \geq 0$ and $k_1 + k_2 \geq 0$.

  By definition, it follows that the connective part of the induced $t$-structure on $\Mod(\Sp_{C_2, i2}; \uZt) \otimes_{\Z} \IndCoh(\Mfg)$ is generated under colimits and extensions by the tensor products $\Sigma^{k_1 + k_2 \sigma} \uZt \otimes \omega^{\otimes n}_{\mathbb{G} / \Mfg}$, where $k_1 \geq 0$ and $k_1 + k_2 \geq 0$. Under the equivalence of \Cref{thm:Cta}, these correspond to $\Sigma^{n+k_1, n+k_2, n} C\ta$, as desired. \todo{Check sign of $n$.}

  It follows directly from the above identification of the tensor connective category that a $C\ta$-module is is tensor coconnective if and only if it is Chow coconnective.
  By Chow connectivity of $C\ta$, which follows from \Cref{lem:Cta-pi0}, it also follows that a $C\ta$-module is Chow connective if it is tensor connective.
  On the other hand, if $X$ is Chow connective, then $X \otimes C\ta$ is clearly tensor connective. Using the equivalence $X \otimes C\ta \simeq X \oplus \Sigma^{1,0,-1} X$ (since $X$ is a $C\ta$-module), we see that $X$ itself is also tensor connective.

  As a consequence of the above, we obtain a symmetric monoidal functor
  \[ \Mod(\SHRat; C\ta)^{\th} \to \SHRatch. \]
  Since the tensor $t$-structure is induced by the Chow $t$-structure, the fact that $C\ta$ lies in the Chow heart from \Cref{lem:Cta-pi0} implies that the above functor factors through a symmetric monoidal equivalence
  \[ \Mod(\SHRat; C\ta)^{\th} \simeq \Mod(\SHRatch; C\ta). \]
  Finally, the equivalence $\pi_0 ^{\ch} \Ss_2 ^{0,0,0} \simeq C\ta$ of \Cref{lem:Cta-pi0} implies that $C\ta$ is the unit of $\SHRatch$, so that the forgetful functor
  \[\Mod(\SHRatch; C\ta) \to \SHRatch\]
  is a symmetric monoidal equivalence.
\end{proof}

We now move on to the proof of \Cref{thm:pi0-MGL}. First we need to recall some basic facts about cell structures on finite approximations of $\MGL$.

\begin{dfn}
  Let $\mathrm{Gr}_k (\mathbb{A}^{n+k})$ denote the Grassmannian scheme of $k$-planes in $\mathbb{A}^{n+k}$, and let $\gamma$ denote the tautological $k$-dimensional vector bundle over $\mathrm{Gr}_k (\mathbb{A}^{n+k})$.

  We let $\MGL(n,k)$ denote the Thom spectrum of the virtual bundle $\gamma - k\cdot \mathrm{triv}$. Then there is a canonical map $\MGL (n_1,k_1) \to \MGL(n_2, k_2)$ whenever $n_2 \geq n_1$ and $k_2 \geq n_2$, and we have $\varinjlim \MGL(n,n) \simeq \MGL$.
\end{dfn}

\begin{prop}[\cite{DICell}] \label{prop:MGL-chow-conn} \label{cor:MGL-dual-Chow-conn}
  The $\R$-motivic spectrum $\MGL(n,k)$ admits a finite cell structure with cells of the form $\Ss^{m,m,m}$.

  As a consequence, $\MGL(n,k)_2$, its Spanier-Whitehead dual $\mathbb{D}(\MGL(n,k)_2)$, and $\MGL_2$ are all Chow connective.
\end{prop}

\begin{proof}
  The cell structure is a consequence of \cite{DICell}.
  It follows from \Cref{thm:complete-compare} and \Cref{prop:compl-facts}(1) that there are equivalences:
  \begin{align*}
    &\MGL(n,k) \otimes \Ss_2 \simeq \MGL(n,k)_2 \\
    &\mathbb{D} (\MGL(n,k)) \otimes \Ss_2 \simeq \mathbb{D} (\MGL(n,k)_2) \\
    &\MGL \otimes \Ss_2 \simeq \MGL_2.
  \end{align*}
  Therefore these spectra continue to have the same cell structure after $2$-completion, from which the statements about Chow connectivity follow immediately.
\end{proof}

%\begin{lem}
  %There is an equivalence of $\mathbb{E}_\infty$-rings $\MGL \simeq \Gamma_* (\MUR)$.
%\end{lem} \todo{Push this one into Section 4!}
%
%\begin{proof}
  %Let us view Artin-Tate $\R$-motivic spectra as filtered $\Gamma_* (\Ss)$-modules.
  %First, we note that since there is an extra degeneracy we have $\Gamma_* (\MUR) \simeq \tau_{\geq 0} \MUR$.
%
  %On the other hand, we have ...
%\end{proof}

\begin{prop} \label{prop:MGL-Cta}
  The map $\MGL_2 \to \MGL_2 \otimes C\ta$ induces an equivalence $\pi_0 ^{\ch} \MGL_2 \simeq \MGL_2 \otimes C\ta$, and the corresponding $\underline{(\MU_2)_{2*} \MU_2}$-comodule is $\underline{(\MU_2)_{2*} \MU_2}$.
  %We have $\MGL \otimes C\ta \in \SHRatch$. Under the equivalence $\SHRatch \simeq \Comod(\uAbh; \underline{\MU_{2*} \MU})$, it corresponds to $\MU_{2*}$
\end{prop}

\begin{proof}
  The first part of the statement follows directly from \Cref{cor:cta-homotopy} and \Cref{prop:MGL-chow-conn}.
  On the other hand, \Cref{cor:gamma-MGL} states that $\MGL_2 \simeq \Gamma_* (\MU_{\R,2})$, so that \Cref{thm:Cta-homotopy} and \Cref{prop:t-compare} together imply that $\MGL_2 \otimes C\ta$ coreponds to $\piu_{*\rho} ^{C_2} (\MU_{\R,2} \otimes \MU_{\R,2}) \cong \underline{(\MU_2)_{2*} \MU_2}$.
  %Since $\MGL$ is Chow connective by \Cref{prop:MGL-chow-conn}, the map $\MGL \to \MGL \otimes C\ta$ induces an equivalence on $\pi_0 ^{\ch}$. On the other hand, \Cref{cor:gamma-MGL}, \Cref{thm:Cta-homotopy} and \Cref{lem:tensor-t-ta} together imply that $\MGL \otimes C\ta$ lies in the heart of the Chow $t$-structure and corresponds to $\piu_{*\rho} ^{C_2} (\MUR \otimes \MUR) \cong \underline{\MU_{2*} \MU}$.
\end{proof}

\begin{ntn}
  We let $\otimes^{\heartsuit}$ denote the symmetric monoidal structure on $\SHRatch \simeq \Comod(\uAbh_2; \underline{(\MU_2)_{2*} \MU_2})$.

  Given a graded $\underline{(\MU_2)_{2*} \MU_2}$-comodule $M_*$, we let $M[n]_*$ denote the shift with $M[n]_{k+n} = M_{k}$.
\end{ntn}

\begin{lem} \label{lem:MGL-heart-tensor}
  Let $M_*$ be a graded $\underline{(\MU_2)_{2*} \MU_2}$-comodule in $C_2$-Mackey functors, viewed as a element of $\SHRatch$.
  Then $M_* \otimes \MGL_2$ lies in $\SHRatch$ and corresponds to $M_* \otimes^{\heartsuit} \underline{(\MU_2)_{2*} \MU_2}$.
\end{lem}

\begin{proof}
  It follows from \Cref{prop:MGL-chow-conn} that $M_* \otimes \MGL_2$ is Chow connective.
  As a consequence, \Cref{prop:MGL-Cta} implies that there are equivalences
  \begin{align*}
    \pi_0 ^{\ch} (M_* \otimes \MGL_2) &\simeq M_* \otimes^{\heartsuit} \pi_0 ^{\ch} \MGL_2 \\
    &\simeq M_* \otimes^{\heartsuit} \underline{(\MU_2)_{2*} \MU_2}.
  \end{align*}

  It therefore suffices to show that $M_* \otimes \MGL_2$ is Chow coconnective.
  For this, we compute for $n,k_1, k_2 \in \Z$ with $k_1 > 0$ and $k_1+ k_2 > 0$:
  \begin{align*}
    \pi_{n+k_1, n+k_2, n} ^{\R} (M _* \otimes \MGL_2) &\cong [\Ss_2^{n+k_1, n+k_2, n}, M _* \otimes \MGL_2]_{\SHRat} \\
    &\cong \varinjlim [\Ss_2^{n+k_1, n+k_2, n}, M_* \otimes \MGL(m,m)_2]_{\SHRat} \\
    &\cong \varinjlim [\Ss_2^{n+k_1, n+k_2, n} \otimes \mathbb{D} (\MGL(m,m)_2), M_*]_{\SHRat} \\
    &\cong 0
  \end{align*}
  since $\mathbb{D} (\MGL(m,m)_2)$ is Chow connective by \Cref{prop:MGL-chow-conn}.
  It follows that $M_* \otimes \MGL_2$ is Chow coconnective, as desired.
\end{proof}

\begin{lem} \label{lem:MGL-tensor-thing}
  Suppose that $X$ is Chow connective. Then the map $X \otimes \MGL_2 \to \pi_0 ^{\ch} (X) \otimes \MGL_2$ induces an equivalence on $\piu_{n-k_1,n-k_2,n} ^{\R}$ for all $n,k_1,k_2 \in \Z$ with $k_1 \geq 0$ and $k_1 + k_2 \geq 0$.
\end{lem}

\begin{proof}
  It suffices to show that $\piu_{n-k_1, n-k_2, n} ^{\R} Y \otimes \MGL_2 = 0$ for all Chow $1$-connective $Y$.
  Since this condition is closed under filtered colimits, suspensions and extensions, it is closed under all colimits and extensions. It therefore suffices to assume that $Y = \Ss_2^{m+a_1, m+a_2, m}$ for $m,a_1,a_2 \in \Z$ satisfying $a_1 > 0$ and $a_1 + a_2 > 0$.
  In other words, we must show that $\piu_{n-k_1, n-k_2, n} ^{\R} \MGL_2$ for all $n,k_1, k_2 \in \Z$ now satisfying $k_1 > 0$ and $k_1 + k_2 > 0$.

  By \Cref{prop:P-MGL}, there is an isomorphism $\piu_{n-k_1, n-k_2, n} ^{\R} \MGL_2 \cong \piu_{n-k_1 + (n-k_2)\sigma} ^{C_2} P_{2n} \MU_{\R,2}$. This is zero because $\Phi^{e} (\Ss_2^{n-k_1 + (n-k_2)\sigma}) \simeq \Ss_2^{2n-k_1-k_2}$ is of dimension $< 2n$ and $\Phi^{C_2} (\Ss_2^{n-k_1 + (n-k_2) \sigma}) \simeq \Ss_2^{n-k_1}$ is of dimension $< n$. %\todo{Dunno where to find this in the literature... hope I don't need to write a lemma for it.}
\end{proof}

\begin{proof}[Proof of \Cref{thm:pi0-MGL}]
  %Given a graded $\underline{\MU_{2*} \MU}$-comodule $\underline{M}_*$, we let $\underline{M}[n]_*$ denote the shift with $\underline{M}[n]_{k+n} = \underline{M}_{k}$.
%
  We have
  \begin{align*}
    \pi_{0} ^{\ch} (X)_n &\cong \uHom_{\underline{(\MU_2)_{2*}}-\mathrm{mod}} (\underline{(\MU_2)_{2*}}[n], \pi_0 ^{\ch} (X)) \\
    &\cong \uHom_{\underline{(\MU_2)_{2*} \MU_2}-\mathrm{comod}} (\underline{(\MU_2)_{2*}}, \pi_0 ^{\ch} (X) \otimes^{\heartsuit} \underline{(\MU_2)_{2*} \MU_2}) \\
    &\cong \piu_0 ^{C_2} \Map_{\SHRat} (\Ss_2^{n,n,n}, \pi_0 ^{\ch} (X) \otimes \MGL_2) \\
    &\cong \piu_0 ^{C_2} \Map_{\SHRat} (\Ss_2^{n,n,n}, X \otimes \MGL_2) \\
    &\cong \underline{(\MGL_2)}^{\R}_{n+k,n,n} (X),
  \end{align*}
  where the third and fourth isomorphisms follow from \Cref{lem:MGL-heart-tensor} and \Cref{lem:MGL-tensor-thing}, respectively.
\end{proof}

%
%\begin{ntn}
%  Let $\tau_{\geq 0} ^{\cn}$ denote connectivization with respect to the Chow $t$-structure.
%\end{ntn}

Finally, we record the following proposition for use in \Cref{sec:nu}.

\begin{prop} \label{prop:complete-compare}
  An Artin-Tate $\R$-motivic spectrum $X \in \SHRat$ is left complete with respect to the Chow $t$-structure if and only if it is $\MGL_2$-local.
\end{prop}

The proof will make use of the following lemma:

\begin{lem} \label{lem:conservative-funcs}
  The functors $X \mapsto \piu_{n+k,n,n} ^{\R} X$ are jointly conservative on $\SHRat$.
\end{lem}

\begin{proof}
  Since the functors $X \mapsto \Map_{\SHRat} (\Ss_2^{n,n,n}, X)$ are jointly conservative, it suffices to note that the functors $Y \mapsto \piu_{k} ^{C_2} Y$ are jointly conservative on $\Sp_{C_2}$.
\end{proof}

\begin{proof}[Proof of \Cref{prop:complete-compare}]
  Since the Chow $t$-structure is right complete, it is equivalent to show that $X \otimes \MGL_2 \simeq 0$ if and only if $\pi_{n} ^{\ch} X = 0$ for all $n$. This is an immediate consequence of \Cref{thm:pi0-MGL} and \Cref{lem:conservative-funcs}.
\end{proof}




\subsection{The functor $\nur$} \label{sec:nu}\ 

In this section, we will study the functor defined below:

\begin{dfn}
  We define a lax symmetric monoidal functor
  \[\nur : \Sp_{C_2, i2} \to \SHRat \]
  to be the composite
  \[\Sp_{C_2, i2} \xrightarrow{\b^{-1}} \Mod(\SHRat; \Ss_2[\ta^{-1}]) \xrightarrow{\tau_{\geq 0} ^{\cn}} \SHRat.\]
\end{dfn}

One of the main results of this section is the following theorem, which summarizes the basic properties of $\nur$:

\begin{thm} \label{thm:nu-properties}
  Let $X, Y, Z \in \Sp_{C_2, i2}$. The lax symmetric monoidal functor $\nur : \Sp_{C_2, i2} \to \SHRat$ satisfies the following properties:
  \begin{enumerate}
    \item There is a natural equivalence $\b(\nur(X)) \simeq X$.
    \item The functor $X \mapsto \nur(X)$ commutes with filtered colimits.
    \item Given any $k \in \Z$, there is a natural equivalence $\nur(\Sigma^{k\rho} X) \simeq \Sigma^{k,k,k} \nur(X)$.
    \item Suppose that $X \to Y \to Z$ is a cofiber sequence. Then
      \[\nur (X) \to \nur (Y) \to \nur (Z)\]
      is a cofiber sequence if and only if $\piu_{n\rho-1} ^{C_2} (X \otimes \MU_{\R, 2}) \to \piu_{n\rho-1} ^{C_2} (Y \otimes \MU_{\R, 2})$ is a monomorphism for all $n$.
    \item Suppose that $X$ is a filtered colimit of $C_2$-spectra that admit finite cell structures with all cells of the form $\Ss_2^{n\rho}$ for $n \in \Z$. Then, for all $Y$, the natural map
      \[\nur (X) \otimes \nur(Y) \to \nur(X \otimes Y)\]
      is an equivalence.
    \item The $C_2$-spectrum $X$ is $\MU_{\R, 2}$-local if and only if $\nur (X)$ is $\MGL_2$-local. 
  \end{enumerate}
\end{thm}

We also compare $\nur$ to the functor $\Gamma_* : \Sp_{C_2, i2} \to \SHRat$ constructed in \Cref{cnstr:gamma-star}.

\begin{thm} \label{thm:Gamma-nu-comp}
  For $X \in \Sp_{C_2, i2}$, there is a natural equivalence
  \[\Gamma_* (X) \simeq \nur (X) ^{\wedge} _{\MGL_2}.\]
\end{thm}

As evidence that the functor $\nur$ is a good way to produce Artin-Tate $\R$-motivic spectra, we prove that several $\R$-motivic spectra of interest may be recovered from their Betti realizations via $\nur$:

\begin{thm}\label{thm:recover}
  There are natural equivalences of commutative rings in $\SHRat$:
  \begin{align*}
    \MGL_2 &\simeq \nur (\MU_{\R,2}) \\
    \HFt &\simeq \nur (\uFt) \\
    \HZt &\simeq \nur (\uZt) \\
    \kgl_2 &\simeq \nur (\ku_{\R,2}) \\
    \kq_2 &\simeq \nur (\ko_{C_2, 2}).
  \end{align*}
\end{thm}

\begin{rmk}
  In light of \Cref{thm:recover}, it would be reasonable to define an $\R$-motivic spectrum of motivic modular forms as $\nur (\tmf_{C_2,2})$, assuming that one had a suitable $C_2$-equivariant spectrum of connective topological modular forms $\tmf_{C_2,2}$.
  Since no such equivariant spectrum has yet been constructed, we leave this to future work. \todo{Mention why previous attempts don't work? Errors in the literature? And perhaps mention that $\TMF_{C_2}$ is known to exist?}
\end{rmk}

\begin{rmk} \label{rmk:gamma-nu-same}
  Since all of the Artin-Tate $\R$-motivic spectra in question are $\MGL_2$-complete, in light of \Cref{thm:Gamma-nu-comp} we may replace $\nur$ in \Cref{thm:recover} by $\Gamma_*$.
\end{rmk}

We now begin with the proof of \Cref{thm:nu-properties}.

\begin{proof}[Proof of \Cref{thm:nu-properties}] \

  \textbf{Proof of (1):} Equivalently, we are to show that $\nur(X)[\ta^{-1}] \simeq \b^{-1} (X)$.
  Using the cofiber sequence $\nur(X) \to \b^{-1} (X) \to \tau_{\leq -1} ^{\cn} (\b^{-1} (X))$, we find that the cofiber of $\nur(X) [\ta^{-1}] \to \b^{-1} (X)[\ta^{-1}] \simeq \b^{-1} (X)$ is may be expressed as $\varinjlim \Sigma^{0,0,m} \tau_{\leq -1} ^{\cn} (\b^{-1} (X))$.

  Now, we know that $\pi_{n+k_1, n+k_2, n} ^{\R} (\tau_{\leq -1} ^{\cn} (\b^{-1} (X))) \cong 0$ for all $n,k_1,k_2 \in \Z$ with $k_1 > 0$ and $k_1+k_2 > 0$.
  It follows that $\pi_{n+k_1, n+k_2, n} ^{\R} ( \Sigma^{0,0,m} \tau_{\leq -1} ^{\cn} (\b^{-1} (X))) \cong 0$ for all $n,k_1,k_2 \in \Z$ with $k_1 > -m$ and $k_1+k_2 > -2m$.
  This implies that $\varinjlim \Sigma^{0,0,m} \tau_{\leq -1} ^{\cn} (\b^{-1} (X))$ is null, as desired.

  \textbf{Proof of (2):} This is immediate from the fact that the Chow $t$-structure is compatible with filtered colimits.

  \textbf{Proof of (3):} It is clear that there is a natural equivalence $\b^{-1} (\Sigma^{k\rho} X) \simeq \Sigma^{k,k,k} \b^{-1} (X)$.
  Moreover, since $\Sigma^{m,m,m}$ induces an automorphism of the Chow $t$-structure, we find further that $\tau_{\geq 0} ^{\cn}$ commutes with $\Sigma^{m,m,m}$.
  Combining these facts, we obtain the desired result.

  \textbf{Proof of (4):} It is clear that $X \to Y \to Z$ is a cofiber sequence if and only if $\b^{-1} (X) \to \b^{-1} (Y) \to \b^{-1} (Z)$ is.
  Applying a general fact about $t$-strutures, we see that
  \[\tau_{\geq 0} ^\cn (\b^{-1} (X)) \to \tau_{\geq 0} ^\cn (\b^{-1} (Y)) \to \tau_{\geq 0} ^\cn (\b^{-1} (Z)) \]
  is a cofiber sequence if and only if $\pi_{-1} ^{\ch} (\b^{-1} (X)) \to \pi_{-1} ^{\ch} (\b^{-1} (Y))$ is a monomorphism.
  The result then follows from \Cref{thm:pi0-MGL} and the isomorphism $\piu_{p,q,w} ^\R (\MGL_2 \otimes \b^{-1} (X)) \cong \piu_{p+q\sigma} ^{C_2} (\MU_{\R, 2} \otimes X)$.

  \textbf{Proof of (5):} By compatibility with filtered colimits, we may assume that $X$ admits a finite cell structure with all cells of the form $\Ss_2^{n\rho}$. Let us write
  \[\Ss_2^{n_1 \rho} \simeq X_1 \to X_2 \to \dots \to X_k = X\]
  with $X_i / X_{i-1} \simeq \Ss_2^{n_i \rho}$.
  Then it is easy to prove that $\piu_{n\rho-1} ^{C_2} (X_i \otimes \MU_{\R, 2}) = 0$ for all $n$ and $i$ by induction, so that by (3) and (4) there are cofiber sequences
  \[\nur (X_{i-1}) \to \nur (X_i) \to \Ss_2^{n_i, n_i, n_i}.\]

  It follows that $\mathbb{D}(\nur(X))$ is Chow connective. To show that the natural map $\nur(X) \otimes \nur(Y) \to \nur(X \otimes Y)$ is an equivalence, it suffices to show that it induces an equivalence on $\Omega^{\infty} \Map_{\SHRat} (Z, -)$ for all Chow connective $Z$.
  For such $Z$, we have
  \begin{align*}
    \Omega^{\infty} \Map_{\SHRat} (Z, \nur(X) \otimes \nur(Y))
    &\simeq  \Omega^{\infty} \Map_{\SHRat} (Z \otimes \mathbb{D} (\nur(X)), \nur(Y)) \\
    &\simeq \Omega^{\infty} \Map_{\SHRat} (Z \otimes \mathbb{D} (\nur(X)), \b^{-1} (Y)) \\
    &\simeq \Omega^{\infty} \Map_{\SHRat} (Z, \nur(X) \otimes \b^{-1} (Y)) \\
    &\simeq \Omega^{\infty} \Map_{\SHRat} (Z, \b^{-1} (X) \otimes \b^{-1} (Y)) \\
    &\simeq \Omega^{\infty} \Map_{\SHRat} (Z, \nur(X \otimes Y)),
  \end{align*}
  as desired. We have used Chow connectivity of $\mathbb{D} (\nur(X))$ in the second equivalence and (1) in the fourth.

  \textbf{Proof of (6):} It is clear that $X$ is $\MU_{\R, 2}$-local if and only $\b^{-1} (X)$ is $\MGL_2$-local.
  By \Cref{prop:complete-compare}, it suffices to show that $\b^{-1} (X)$ is left complete with respect to the Chow $t$-structure if and only if $\nur(X) = \tau_{\geq 0} ^{\cn} (\b^{-1} (X))$ is.
  This follows from the fact that left completeness is invariant under taking connective cover.
\end{proof}

We now move on to the proof of \Cref{thm:Gamma-nu-comp}.
We begin by studying the interaction of the Chow $t$-structure with $\MGL_2$.

\begin{prop} \label{prop:Chow-MGL}
  Let $X \in \SHRat$. Then the following statements hold:
  \begin{enumerate}
    \item If $X$ is Chow connective, then
      \[\Map_{\SHRat} (\Ss_2^{0,0,n}, \MGL_2 \otimes X)\]
      is regular slice $2n$-connective for all $n$.
    \item There is a natural equivalence
      \[\MGL_2 \otimes (\tau^{\cn} _{\geq 0} X) \simeq \tau_{\geq 0} ^{\cn} (\MGL_2 \otimes X).\]
    \item The natural map
      \[\Map_{\SHRat} (\Ss_2^{0,0,n}, \MGL_2 \otimes \tau_{\geq 0} ^{\cn} X) \to \Map_{\SHRat} (\Ss_2^{0,0,n}, \MGL_2 \otimes X)\]
      induced by the counit $\tau_{\geq 0} ^{\cn} X \to X$ is equivalent to the canonical map
      \[P_{2n} \Map_{\SHRat} (\Ss_2^{0,0,n}, \MGL_2 \otimes X) \to \Map_{\SHRat} (\Ss_2^{0,0,n}, \MGL_2 \otimes X).\]
  \end{enumerate}
\end{prop}

\begin{proof}
  We begin with the proof of (1).
  Since the full subcategory of $X$ for which $\Map_{\SHRat} (\Ss_2^{0,0,n}, \MGL_2 \otimes X)$ is regular slice $2n$-connective is closed under colimits and extensions, it suffices to prove this for $X = \Ss_2^{m+k_1,m+k_2,m}$ for all $m,k_1,k_2 \in \Z$ where $k_1 \geq 0$ and $k_1 + k_2 \geq 0$.

  %In the case that $X = \Spec \C$, we have $\pi_{p,q,w} ^{\R} (\MGL_2 \otimes \Spec \C) \cong \pi_{p+q,w} ^{\C} \MGL_2$.
  %This is $0$ when $p+q < w$, from which the result follows.
%
  Now, we have
  \[\Map_{\SHRat} (\Ss_2^{0,0,n}, \MGL_2 \otimes \Ss_2^{m+k_1,m+k_2,m}) \simeq \Sigma^{m+k_1+(m+k_2)\sigma} \Map_{\SHRat} (\Ss_2^{0,0,n-m}, \MGL_2).\]
  Since $\Sigma^{m+k_1+(m+k_2)\sigma}$ of a regular slice $2(n-m)$-connective $C_2$-spectrum is regular slice $2n$-connective, this reduces us to the case where $X = \Ss_2^{0,0,0}$, where this follows from \Cref{prop:P-MGL}.

  We now prove (2).
  Since $\MGL_2$ is Chow connective, there is a natural map 
  \[(\tau_{\geq 0} ^{\cn} X) \otimes \MGL_2 \to \tau_{\geq 0} ^{\cn} (X \otimes \MGL_2).\]
  To show that this is an equivalence, it suffices to show that
  \[\Omega^{\infty} \Map_{\SHRat} (Y, (\tau_{\geq 0} ^{\cn} X) \otimes \MGL_2) \to \Omega^{\infty} \Map_{\SHRat} (Y, \tau_{\geq 0} ^{\cn} (X \otimes \MGL_2))\]
  is an equivalence for all $Y$ compact and Chow connective.
  This is equivalent to the map
  \[\Omega^{\infty} \Map_{\SHRat} (Y, (\tau_{\geq 0} ^{\cn} X) \otimes \MGL_2) \to \Omega^{\infty} \Map_{\SHRat} (Y, X \otimes \MGL_2).\]
  Now, using compactness of $Y$, the identification $\MGL_2 \simeq \varinjlim_{k} \MGL(k,k)_2$, and duality, we may rewrite this map as
  \[\varinjlim_k \Omega^{\infty} \Map_{\SHRat} (Y \otimes \mathbb{D} (\MGL(k,k)_2), (\tau_{\geq 0} ^{\cn} X)) \to \varinjlim_k \Omega^{\infty} \Map_{\SHRat} (Y \otimes \mathbb{D} (\MGL(k,k)_2), X).\]
  This is an equivalence by \Cref{cor:MGL-dual-Chow-conn}.

  Finally, we prove (3).
  By (1), the map
  \[\Map_{\SHRat} (\Ss_2^{0,0,n}, \MGL_2 \otimes \tau_{\geq 0} ^{\cn} X) \to \Map_{\SHRat} (\Ss_2^{0,0,n}, \MGL_2 \otimes X)\]
  factors through a map
  \[\Map_{\SHRat} (\Ss_2^{0,0,n}, \MGL_2 \otimes \tau_{\geq 0} ^{\cn} X) \to P_{2n} \Map_{\SHRat} (\Ss_2^{0,0,n}, \MGL_2 \otimes X).\]

  Since both sides are regular slice $2n$-connective, it suffices to show that this map is an equivalence after applying $\Omega^{\infty} \Sigma^{-n\rho}$. Making some basic manipulations, this is equivalent to showing that
  \[\Omega^{\infty} \Map_{\SHRat} (\Ss_2^{n,n,n}, \MGL_2 \otimes \tau_{\geq 0} ^{\cn} X) \to \Omega^{\infty} \Map_{\SHRat} (\Ss_2^{n,n,n}, \MGL_2 \otimes X)\]
  is an equivalence.
  This is equivalence by (2) and the Chow connectivity of $\Ss_2^{n,n,n}$.
\end{proof}
%
%\begin{prop} \label{prop:Chow-conn-slice}
%  Suppose that $X$ is Chow connective. Then \[\Map_{\SHRat} (\Ss^{0,0,n}, \MGL \otimes X)\] is $2n$-slice connective for all $n$.
%\end{prop}
%
%\begin{proof}
%  Since the full subcategory of $X$ for which $\Map_{\SHRat} (\Ss^{0,0,n}, \MGL \otimes X)$ is $2n$-slice connective is closed under colimits, it suffices to prove this for $X = \Ss^{m,m,m}$ for all $m \in \Z$ and $X = \Spec \C$.
%
%  In the case that $X = \Spec \C$, we have $\pi_{p,q,w} ^{\R} (\MGL \otimes \Spec \C) \cong \pi_{p+q,w} ^{\C} \MGL$.
%  This is $0$ when $p+q < w$, from which the result follows.
%
%  In the case that $X = \Ss^{m,m,m}$, we have
%  \[\Map_{\SHRat} (\Ss^{0,0,n}, \MGL \otimes \Ss^{m,m,m}) \simeq \Sigma^{m\rho} \Map_{\SHRat} (\Ss^{0,0,n-m}, \MGL).\]
%  Since $\Sigma^{m\rho}$ of a slice $2(n-m)$-connective $C_2$-spectrum is slice $2n$-connective, this reduces us to the case where $X = \Ss^{0,0,0}$, where this follows from REF.
%\end{proof}
%
%\begin{prop} \label{prop:MGL-Chow-conn-comm}
%  Suppose that $X \in \SHR^{\at}$. Then there is a natural equivalence
%  \[(\tau_{\geq 0} ^{\cn} X) \otimes \MGL \simeq \tau_{\geq 0} ^{\cn} (X \otimes \MGL).\]
%\end{prop}
%
%\begin{proof}
%  Since $\MGL$ is Chow connective, there is a natural map 
%  \[(\tau_{\geq 0} ^{\cn} X) \otimes \MGL \to \tau_{\geq 0} ^{\cn} (X \otimes \MGL).\]
%  Since both sides are Chow connective, to show that this is an equivalence, it suffices to show that
%  \[\Omega^{\infty} \Map_{\SHRat} (Y, (\tau_{\geq 0} ^{\cn} X) \otimes \MGL) \to \Omega^{\infty} \Map_{\SHRat} (Y, \tau_{\geq 0} ^{\cn} (X \otimes \MGL))\]
%  is an equivalence for $Y = \Spec \C$ or $\Ss^{m,m,m}$.
%  This is equivalent to the map
%  \[\Omega^{\infty} \Map_{\SHRat} (Y, (\tau_{\geq 0} ^{\cn} X) \otimes \MGL) \to \Omega^{\infty} \Map_{\SHRat} (Y, X \otimes \MGL).\]
%  Now, using compactness of $Y$, the identification $\MGL \simeq \varinjlim_{k} \MGL(k,k)$, and duality, we may rewrite this map as
%  \[\varinjlim_k \Omega^{\infty} \Map_{\SHRat} (Y \otimes \mathbb{D} (\MGL(k,k)), (\tau_{\geq 0} ^{\cn} X)) \to \varinjlim_k \Omega^{\infty} \Map_{\SHRat} (Y \otimes \mathbb{D} (\MGL(k,k)), X).\]
%  By Chow connectivity of $\mathbb{D} (\MGL(k,k))$, this is an equivalence.
%\end{proof}
%
%\begin{prop} \label{prop:MGL-Chow}
%  Let $X$ be an $\R$-motivic spectrum. The map
%  \[\Map_{\SHRat} (\Ss^{0,0,n}, \MGL \otimes \tau_{\geq 0} ^{\cn} X) \to \Map_{\SHRat} (\Ss^{0,0,n}, \MGL \otimes X)\]
%  induced by the (co?)unit $\tau_{\geq 0} ^{\cn} X \to X$ is equivalent to the canonical map
%  \[P_{2n} \Map_{\SHRat} (\Ss^{0,0,n}, \MGL \otimes X) \to \Map_{\SHRat} (\Ss^{0,0,n}, \MGL \otimes X).\]
%\end{prop}
%
%\begin{proof}
%  By \Cref{prop:Chow-conn-slice}, the map
%  \[\Map_{\SHRat} (\Ss^{0,0,n}, \MGL \otimes \tau_{\geq 0} ^{\cn} X) \to \Map_{\SHRat} (\Ss^{0,0,n}, \MGL \otimes X)\]
%  factors through a map
%  \[\Map_{\SHRat} (\Ss^{0,0,n}, \MGL \otimes \tau_{\geq 0} ^{\cn} X) \to P_{2n} \Map_{\SHRat} (\Ss^{0,0,n}, \MGL \otimes X).\]
%
%  Since both sides are slice $2n$-connective, it suffices to show that this map is an equivalence after applying $\Omega^{\infty} \Sigma^{-n\rho}$. Making some basic manipulations, this is equivalent to showing that
%  \[\Omega^{\infty} \Map_{\SHRat} (\Ss^{n,n,n}, \MGL \otimes \tau_{\geq 0} ^{\cn} X) \to \Omega^{\infty} \Map_{\SHRat} (\Ss^{n,n,n}, \MGL \otimes X)\]
%  is an equivalence.
%
%  This is equivalence by \Cref{prop:MGL-Chow-conn-comm} and the Chow connectivity of $\Ss^{n,n,n}$.
%%
  %To show this, we use the identification $\MGL \simeq \varinjlim_k \MGL(k,k)$, compactness of $\Ss^{n,n,n}$, and duality to rewrite this as
  %\[\varinjlim_k \Omega^{\infty} \Map_{\SHRat} (\Ss^{n,n,n} \otimes \mathbb{D}(\MGL(k,k)), \tau_{\geq 0} ^{\cn} X) \to \varinjlim_k \Omega^{\infty} \Map_{\SHRat} (\Ss^{n,n,n} \otimes \mathbb{D}(\MGL(k,k)), X).\]
  %\todo{Need to justify why I can relate $t$-structure to underlying $C_2$-space instead of just space.}
%
  %This map is an equivalence since $\Ss^{n,n,n} \otimes \mathbb{D} (\MGL(k,k))$ is Chow connective by \Cref{cor:MGL-dual-Chow-conn}.
%\end{proof}

We are now ready to prove \Cref{thm:Gamma-nu-comp}.
%
%\begin{prop}
%  Given $X \in \SHR^{\at}$, there are natural equivalences
%  \[\Gamma_* (\b(X)) \simeq (\tau_{\geq 0} ^{\cn} (X[\ta^{-1}]))^{\wedge} _{\MGL},\]
%  where the $\MGL$-completion is taken internally to $\SHR^{\at}$.
  %\simeq \tau_{\geq 0} ^{\cn} ((X[\ta^{-1}])^{\wedge} _{\MGL}).\]
  %
  %
  %In particular, if $\b(X)$ is $\MUR$-complete, then there is a natural equivalence
  %
  %\[\Gamma_* (\b(X)) \simeq \tau_{\geq 0} ^{\cn} (X[\ta^{-1}]).\]
%\end{prop}
%
\begin{proof}[Proof of \Cref{thm:Gamma-nu-comp}]
  In this proof, we freely use the notation of \Cref{sec:top} and \Cref{app:compare}, in particular the functors $i_*$, $\tau_{\geq 0} ^{2\slice}$, and $\cb$.

  Let $X$ be a $C_2$-spectrum. By \Cref{prop:Chow-MGL}(2), there are equivalences
  \begin{align*}
    \nur(X)^{\wedge} _{\MGL_2}
    &\simeq \Tot^{\bullet} (\nur(X) \otimes \cb(\MGL_2)) \\
    &\simeq \Tot^{\bullet} (\nur(X \otimes \cb(\MU_{\R, 2}))).
  \end{align*}
  Applying $i_*$ to the map $\nur(X \otimes \cb(\MU_{\R, 2})) \to \b^{-1} (X \otimes \cb(\MU_{\R, 2}))$ of $\cb(\MGL_2)$-modules in cosimplicial Artin-Tate $\R$-motivic spectra, we obtain a map
  \[i_* (\nur (X \otimes \cb(\MU_{\R, 2}))) \to X \otimes \cb(\MU_{\R, 2})\]
  of $i_* (\cb(\MGL_2)) \simeq \tau_{\geq 0}^{2\slice} \cb(\MU_{\R, 2})$-modules in cosimplicial filtered $C_2$-spectra, where the target is constant in the filtered direction.

  By \Cref{prop:Chow-MGL}(3), this factors through an equivalence
  \[i_* (\nur (X \otimes \cb(\MU_{\R, 2}))) \xrightarrow{\sim} \tau_{\geq 0} ^{2\slice} (X \otimes \cb(\MU_{\R, 2}))\]
  of $\tau_{\geq 0} ^{2\slice} \cb(\MU_{\R, 2})$-modules.
  Totalizing, we obtain an equivalence
  \[i_* (\nur (X)^{\wedge} _{\MGL_2}) \simeq i_* (\Gamma_* (X))\]
  of $R_{\bullet}$-modules.
  Finally, translating back along the equivalence $i_*$, we obtain the desired equivalence
  \[\nur(X)^{\wedge} _{\MGL_2} \simeq \Gamma_* (X)\]
  in $\SHRat$.
%
%  In this proof, we view $\SHR^{\at}$ as the category of filtered $\Gamma_* (\Ss)$-modules under the equivalence of \Cref{thm:filt-model}.
%  By definition, $\nur(X)$ is $\tau_{\geq 0} ^{\cn}$ applied to the $\Gamma_* (\Ss)$-module
  %Let us identify $X$ with the $\Gamma_* (\Ss)$-module
  %\[\dots \to i_* (X)_1 \to i_* (X)_0 \to i_* (X)_{-1} \to \dots.\]
  %Then we may identify $X[\ta^{-1}]$ with
%  \[\dots \xrightarrow{\mathrm{id}} X \xrightarrow{\mathrm{id}} X \xrightarrow{\mathrm{id}} X \xrightarrow{\mathrm{id}} \dots.\]
%  Applying \Cref{prop:Chow-MGL}(3), we may identify $\nur (X) \otimes \MGL$ with
%  \[ \dots \to P_2 (X \otimes \MUR) \to P_0 (X \otimes \MUR) \to P_{-2} (X \otimes \MUR) \to \dots. \]
%  Using the splitting of $\MGL \otimes \MGL$, we infer more generally that the cosimplical object $\nur (X) \otimes \MGL^{\bullet + 1}$ is equivalent to
%  \[ \dots \to P_2 (X \otimes \MUR^{\bullet + 1}) \to P_0 (X \otimes \MUR^{\bullet+1}) \to P_{-2} (X \otimes \MUR^{\bullet+1}) \to \dots. \]
%  Totalizing, we find that 
%  \[\nur(X)^{\wedge} _{\MGL} \simeq \Tot^\bullet (\nur(X) \otimes \MGL^{\bullet+1}) \simeq \Tot^{\bullet} (P_{2*} (X \otimes \MUR^{\bullet+1})) \simeq \Gamma_* (X),\]
%  as desired.
\end{proof}

%In the second half of this section, we prove that several $\R$-motivic spectra of interest can be recovered from their Betti realizations via the functor $\nur$. \todo{As of now, there is a slight gap in the proof, really I show that you get the $\MGL$-localization of $nu$ of stuff, so you either need to replace $\nur$ with $\Gamma_*$ here or prove that sometimes $\nur$ of stuff is $\MGL$-local.}
%
%\begin{thm}\label{thm:recover}
  %There are natural equivalences in $\SHR_{i2}$:
  %\[\MGL_2 \simeq \nur (\MU_{\R,2})\]
  %\[\HFt \simeq \nur (\uFt)\]
  %\[\HZt \simeq \nur (\uZt)\]
  %\[\kgl_2 \simeq \nur (\ku_{\R,2})\]
  %\[\kq_2 \simeq \nur (\ko_{C_2, 2}).\]
%\end{thm}
%
%\begin{rmk}\label{rmk:nu-gamma}
  %Since all of the $\R$-motivic spectra appearing in \Cref{thm:recover} are $2$-complete and connective for the homotopy $t$-structue, they are $\MGL$-complete. It follows from \Cref{thm:Gamma-nu-comp} that we could replace $\nur$ by $\Gamma_*$ in \Cref{thm:recover}.
%\end{rmk}
We now move on to the proof of \Cref{thm:recover}.
We will prove \Cref{thm:recover} as an application of a general criterion for there to be an $\MGL_2$-local equivalence $X \simeq \nur (\b(X))$. This will be based on the following definition:

%\begin{dfn}
  %We say that $X \in \SHRat$ is \emph{slice simple} if the following conditions hold:
 %\begin{enumerate}
   %\item The groups $\pi_{k,k,*} ^\R X$ and $\pi_{k+1, k, *}^{\R} X /\ker(\res)$ are $\ta$-torsion free for all $k \in \Z$, where $\res : \pi_{p,q,w} ^\R \to \pi_{p+q,w} ^\C$ is the map induced by base change, and the groups $\pi_{k,*} ^\C X_{\C}$ are $\tau$-torsion free for all $k \in \Z$.
   %\item The groups $\pi_{k,k,n} ^\R X$ and $\pi_{k+1,k,n}^\R X/\ker(\res)$ are zero for $k < n$, and $\pi_{k,n} ^{\C} X_{\C}$ is zero for $k < 2n$.
    %\item The groups $\pi_{k,k,n} ^\R X \otimes C\ta$ and $\pi_{k+1,k,n} ^\R X \otimes C\ta / \ker(\res)$ are zero for $k \neq n$, and the groups $\pi_{k,n} ^{\C} X_{\C} \otimes C\tau$ are zero for $k \neq 2n+1, 2n$.
  %\end{enumerate}
%\end{dfn}

\begin{dfn}
  We say that $X \in \SHRat$ is \emph{slice simple} if, for all $n \in \Z$, the natural map
  \[ \Map_{\SHRat} (\Ss_2^{0,0,n}, X) \to \b(X)\]
  factors through an equivalence
  \[ \Map_{\SHRat} (\Ss_2^{0,0,n}, X) \simeq P_{2n} \b(X).\]
\end{dfn}
%
%\begin{lem} \label{lem:slice-simple-alt}
%  Let $X \in \SHR^{\at}$. Then $X$ is slice simple if and only if the natural map
%  \[ \Map_{\SHRat} (\Ss^{0,0,n}, X) \to \b(X)\]
%  factors through an equivalence
%  \[ \Map_{\SHRat} (\Ss^{0,0,n}, X) \simeq P_{2n} \b(X).\]
%\end{lem}
%
%\begin{proof}
%  Let us begin by giving equivalent reformulations of the conditions for a spectrum to be slice simple.
%
%  Condition (1) corresponds to injectivity of the natural maps $\pi_{k,k,n} ^\R X \to \pi_{k\rho} ^{C_2} \b(X)$, $\pi_{k+1,k,n} ^{\R} X/\ker(\res) \to \pi_{k\rho+1} ^{C_2} \b(X) /\ker(\res)$, and $\pi_{k,n} ^{\C} X_{\C} \to \pi_{k} X_{\C}$.
%
%  Condition (2) corresponds to slice $2n$-connectivity of $\Map_{\SHRat} (\Ss^{0,0,n}, X)$.
%
%  Given condition (2), condition (3) corresponds to surjectivity of the natural maps $\pi_{k,k,n} ^\R X \to \pi_{k\rho} ^{C_2} \b(X)$ and $\pi_{k+1,k,n} ^{\R} X/\ker(\res) \to \pi_{k\rho+1} ^{C_2} \b(X) /\ker(\res)$ for $k \geq n$ and $\pi_{k,n} ^{\C} X_{\C} \to \pi_{k} X_{\C}$ for $k \geq 2n$.
%
%  Using the identification of the regular $2n$-slices of a $C_2$-spectrum $Y$ with $\Sigma^{n\rho} \piu ^{C_2} _{n\rho} Y$ and the regular $2n+1$-slices with $\Sigma^{n\rho+1} \piu ^{C_2} _{n\rho+1} Y / \ker(\res)$, these conditions are manifestly equivalent to the natural map
%  \[ \Map_{\SHRat} (\Ss^{0,0,n}, X) \to \b(X)\]
%  factoring through an equivalence
%  \[ \Map_{\SHRat} (\Ss^{0,0,n}, X) \simeq P_{2n} \b(X). \qedhere\]
%
%  Suppose that $X$ is slice simple.
%  It follows from condition (2) in the definition of slice simple that $\Map_{\SHRat} (\Ss^{0,0,n}, X)$ is slice $2n$-connective, so that the map
%  \[ \Map_{\SHRat} (\Ss^{0,0,n}, X) \to \b(X)\]
%  factors naturally through a map
%  \[ \Map_{\SHRat} (\Ss^{0,0,n}, X) \to P_{2n} \b(X).\]
%
%  Combining conditions (1), (2) and (3), we learn that the natural maps $\pi_{k,k,n} ^\R X \to \pi_{k\rho} ^{C_2} \b(X)$, $\pi_{k+1,k,n} ^{\R} X/\ker(\res) \to \pi_{k\rho+1} ^{C_2} \b(X) /\ker(\res)$ are isomorphisms for all $k \geq n$, and that $\pi_{k,n} ^{\C} X_{\C} \to \pi_{k} X_{\C}$ are isomorphisms for all $k \geq 2n$.
%
%  We conclude that 
%
%  On the other hand, suppose that the natural map
%  \[ \Map_{\SHRat} (\Ss^{0,0,n}, X) \to \b(X)\]
%  factors through an equivalence
%  \[ \Map_{\SHRat} (\Ss^{0,0,n}, X) \simeq P_{2n} \b(X).\]
%  Condition (2) follows from the slice $2n$-connectivity of $\Map_{\SHRat} (\Ss^{0,0,n}, X)$.
%  Conditions (1) and (3) follow from the fact that
%\end{proof}

\begin{prop} \label{prop:slice-simple-nu}
  Let $X \in \SHR^{\at}$, and let $L_{\MGL_2}$ denote $\MGL_2$-localization. Then $L_{\MGL_2} X \simeq L_{\MGL_2} \nur (\b(X))$ if and only if $\MGL_2 \otimes X$ is slice simple.

  If $X$ is a commutative algebra, then the equivalence $L_{\MGL_2} X \simeq L_{\MGL_2} \nur (\b(X))$ is one of commutative algebras.
%
%
%  under the following assumptions:
%  \begin{itemize}
%    \item The groups $\MGL_{k,k,*} ^\R X$ and $\MGL_{k,*} ^\C X_{\C}$ are $\ta$-torsion free for all $k \in \Z$.
%    \item The groups $\MGL_{k,k,n} ^\R X \otimes C(\ta)$ are zero for $n \neq k$, and the groups $\MGL_{k,n} ^{\C} X_{\C} \otimes C(\tau)$ are zero for $n \neq 2k+1, 2k$.
%      Furthermore, the groups $\MGL_{k+1,k,n} ^\R X \otimes C(\ta) / \ker(\res)$ are zero for $n \neq k$.
%  \end{itemize}
\end{prop}

\begin{proof}
  Suppose that $L_{\MGL_2} X \simeq L_{\MGL_2} \nur(\b(X))$, so that $\MGL_2 \otimes X \simeq \MGL_2 \otimes \nur(\b(X))$.
  Then we have
  \[\Map_{\SHRat} (\Ss_2^{0,0,n}, \MGL_2 \otimes X) \simeq \Map_{\SHRat} (\Ss_2^{0,0,n}, \MGL_2 \otimes \nur(\b(X))\]
  as $C_2$-spectra over $\b(X)$.
  By \Cref{prop:Chow-MGL}(3), the map
  \[\Map_{\SHRat} (\Ss_2^{0,0,n}, \MGL_2 \otimes \nur(\b(X)) \to \b(\MGL_2 \otimes X)\]
  factors through an equivalence
  \[\Map_{\SHRat} (\Ss_2^{0,0,n}, \MGL_2 \otimes \nur(\b(X)) \simeq P_{2n} \b(\MGL_2 \otimes X),\]
  so that $\MGL_2 \otimes X$ is slice simple.

  Now suppose that $\MGL_2 \otimes X$ is slice simple, and
%
%  By definition of $\nur$, there is an equivalence $\nur (\b(X)) \simeq \tau_{\geq 0} ^{\cn} (X[\ta^{-1}])$, and we will work with the latter for the rest of the proof.
%
  consider the following diagram:
  \begin{center}
    \begin{tikzcd}
      X \ar[r] & X[\ta^{-1}] & \\
      \tau_{\geq 0} ^{\cn} X \ar[r] \ar[u] & \tau_{\geq 0} ^{\cn} (X[\ta^{-1}]). \ar[u] 
    \end{tikzcd}
  \end{center}
  If $X$ is a commutative algebra, then this is a diagram of commutative algebras.
  Applying $\Map_{\SHRat} (\Ss_2^{0,0,n}, - \otimes \MGL_2)$, we obtain the diagram:
  \begin{center}
    \begin{tikzcd}
      \Map_{\SHRat} (\Ss_2^{0,0,n}, X \otimes \MGL_2) \ar[r] & \b(X) \otimes \MU_{\R, 2} & \\
      \Map_{\SHRat} (\Ss_2^{0,0,n}, (\tau_{\geq 0} ^{\cn} X) \otimes \MGL_2) \ar[r] \ar[u] & \Map_{\SHRat} (\Ss_2^{0,0,n}, (\tau_{\geq 0} ^{\cn} (X[\ta^{-1}])) \otimes \MGL_2). \ar[u] 
    \end{tikzcd}
  \end{center}
  Finally, applying \Cref{prop:Chow-MGL}(3) and slice simplicity of $\MGL_2 \otimes X$, we may identify this diagram with
  \begin{center}
    \begin{tikzcd}
      P_{2n} (\b(X) \otimes \MU_{\R, 2}) \ar[r] & \b(X) \otimes \MU_{\R, 2} & \\
      P_{2n} (\b(X) \otimes \MU_{\R, 2}) \ar[r] \ar[u] & P_{2n} (\b(X) \otimes \MU_{\R, 2}). \ar[u] 
    \end{tikzcd}
  \end{center}
  In conclusion, we find that the maps $\tau_{\geq 0} ^{\cn} X \to X$ and $\tau_{\geq 0} ^{\cn} X \to \tau_{\geq 0} ^{\cn} (X[\ta^{-1}]) \simeq \nur(\b(X))$ are $\MGL_2$-equivalences, as desired.
%
%
%  By \Cref{prop:MGL-Chow} and condition (2) in the defintion of slice simple, we find that
%  \[X \otimes \MGL \to (\tau_{\geq 0} ^{\cn} X) \otimes \MGL\]
%  is an equivalence, so that $\tau_{\geq 0} ^{\cn} X \to X$ is an $\MGL$-local equivalence.
%
%  On the other hand, $X \otimes \MGL \simeq (\tau_{\geq 0} ^{\cn} X) \otimes \MGL \to (\tau_{\geq 0} ^{\cn} (X[\ta^{-1}])) \otimes \MGL$
  %It follows that the map $\tau_{\geq 0} ^{\cn} X \to L_{\MGL} (\tau_{\geq 0} ^{\cn} (X[\ta^{-1}]))$ factors through a map $X \to L_{\MGL} (\tau_{\geq 0} ^{\cn} (X[\ta^{-1}])).$
%
  %By definition of slice simple, we find that $X$ lies in $\SHR^{\at, \cn} _{\geq 0, i2}$, so that the natural map $X \to X[\ta^{-1}]$ factors through a map $X \to \tau_{\geq 0} ^{\cn} (X[\ta^{-1}])$.
%
  %By a PROP and the definition of slice simple, we see that this map induces an equivalence on $\MGL_2$-homology. Since $X$ is $\eta$-complete, it is $\MGL_2$-complete and so this implies that $X \to \tau_{\geq 0} ^{\cn} (X[\ta^{-1}])$ is an equivalence.
%
  %Examining the definitions, we find that $\Gamma_* \b(X) \simeq \tau_{\geq 0} ^{\cn} X[\ta^{-1}]$, so we win.
\end{proof}

To prove \Cref{thm:recover}, we now need to show that the $\R$-motivic spectra in question are slice simple after being tensored with $\MGL_2$. To this end, we record the following simple proposition:
%
%\begin{lem}
%  The functors
%  \[\pi_{k,k,n} ^\R\]
%  \[\pi_{k+1,k,n} ^\R / \ker(\res) \]
%  \[\pi_{k,n} ^\C \]
%  are jointly conservative on $\SHR ^{gc}$.
%\end{lem}
%
%\begin{proof}
%%  By joint conservativity of $\pi_{p,q,w} ^\R$, the functors
%  \[X \mapsto \Map_{\SHR_{i2}} (\Ss^{0,0,w}, X)\]
%  are jointly conservative on $\SHR^{gc} _{i2}$.
%
%  Then the result follows from the joint conservativity of the functors $\piu_{n\rho}$ and $\piu_{n\rho+1} / \ker(\res)$ on $C_2$-spectra. This in turn follows from the identification of the slices of a $C_2$-spectrum.
%\end{proof}
%
%\begin{lem}
%  Suppose that $X \in \Sp_{C_2, ip}$ is an $\MU_{\R,2}$-module. Then $\Gamma_* X$ is slice simple.
  %Then the the groups
  %$\pi_{k,k,n} ^\R \nur X$ and $\pi_{k+1,k,n} ^\R \nur X / \ker(\res)$ are zero for $n > k$ and $\pi_{k,n} ^\C \nur X$ is zero for $n > 2k$.
  %On the other hand, the maps induced by Betti realization
  %$\pi_{k,k,n} ^\R \nur X \to \pi_{k\rho} ^{C_2} X$ and
  %$\pi_{k+1,k,n} ^\R \nur X  \to \pi_{k\rho+1}^{C_2} X /\ker(\res)$
  %are isomorphisms for $n \leq k$, and the map
  %$ $
  %is an isomorphism for $n \leq k$.
%\end{lem}
%
%\begin{proof}
%  Since $X$ is an $\MUR$-module, the cosimplicial diagram $X \otimes \MUR^{\bullet +1}$ admits a contracting homotopy, hence the diagrams $P_{2n} (X \otimes \MUR^{\bullet+1})$ do as well.
%  It follows that $\Tot (P_{2n} (X \otimes \MUR^{\bullet+1})) \simeq P_{2n} X$, so that $\Gamma_{n} (X) \simeq P_{2n} X$.
%
%  Since we have $P_{2n} ^{2n} X \simeq \Sigma^{n\rho} \piu_{n\rho} X$ and $P_{2n+1} ^{2n+1} X \simeq \Sigma^{n\rho+1} \piu_{n\rho+1} X / \ker(\res)$, the result follows.
%\end{proof}
%
%
\begin{prop} \label{prop:sl-si-cond}
  The collection of slice simple Artin-Tate $\R$-motivic spectra is closed under direct sums, retracts, and pure suspensions $\Sigma^{n,n,n}$.
%
  %Moreover, suppose that $X \in \SHRat$ has convergent slice tower: $X \simeq \varprojlim f_n X$.
  %Moreover, any Artin-Tate $\R$-motivic spectrum with convergent slice tower and whose $n$-slice is a direct sum of $\Sigma^{n,n,n} \HFt$ and $\Sigma^{n,n,n} \HZt$ for all $n$ is slice simple.
\end{prop}
%
%\begin{proof}
%  The closure properties of slice simple $\R$-motivic spectra are clear from the definition.
%
%  The second statement follows from the slice spetral sequence and the following facts about the trigraded homotopy groups of $\HZt$ and $\HFt$, which may be read off from \Cref{thm:pthomology} and \Cref{thm:eqpthom}:
%  \begin{itemize}
%    \item $\pi_{k,k,*} ^\R \HZt \cong \pi_{k,k,*} ^\R \HFt \cong 0$ for $k \neq 0$.
%    \item $\pi_{k+1,k,*} ^\R \HZt \cong \pi_{k+1,k,*} ^\R \HFt \cong 0$ for all $k$.
%    \item $\pi_{0,0,*} ^\R \HZt \cong \Zt[\ta]$ and $\pi_{0,0,*} ^\R \HFt \cong \Ft[\ta]$.
%    \item $\pi_{*,*} ^\C \HZt \cong \Zt[\tau]$ and $\pi_{*,*} ^\C \HFt \cong \Ft[\tau]$. \qedhere
%  \end{itemize}
%\end{proof}
%
%\begin{cor} \label{cor:in-im-Gamma}
%  There are equivalences
%  \[\MGL_2 \simeq \nur(\MU_{\R,2})\]
%  \[\HFt \simeq \nur(\uFt)\]
%  \[\HZt \simeq \nur(\uZt)\]
%  \[\kgl_2 \simeq \nur(\ku_2)\]
%  \[\kq_2 \simeq \nur (\ko_2)\]
%\end{cor}
%
\begin{proof}[Proof of \Cref{thm:recover}]
  By \Cref{prop:slice-simple-nu}, it suffices to show that $\MGL_2 \otimes E$ is slice simple for $E=\MGL_2, \HFt, \HZt, \kgl_2, \kq_2$. Indeed, each of these $E$ is $\MGL_2$-local, and $\nur(\b(E))$ is also $\MGL_2$-local by \Cref{thm:nu-properties}(6), since $\b(E)$ is $\MU_{\R, 2}$-local.
  %Indeed, as note in \Cref{rmk:nu-gamma}, they are even $\MGL$-complete.

  Now, for $E=\MGL_2, \HFt, \HZt, \kgl_2$, $\MGL_2 \otimes E$ is a direct sum of pure suspensions of $E$, so it suffices to show that $E$ itself is slice simple.
  For $\MGL_2$, this follows from \Cref{prop:P-MGL}. For $\HZt$ and $\HFt$, it follows from \Cref{lem:Zcalc}.
  For $\kgl_2$, it follows from \Cref{thm:eff-fil} and \cite[Theorem 1.1]{SliceComp}.

  %For $\MGL$, this follows from \Cref{prop:sl-si-cond}, since $\MGL \otimes \MGL$ is a direct sum of pure suspensions of $\MGL$ and the slices of $\MGL$ are of the specified form.
  %Since $\HFt$, $\HZt$, and $\kgl_2$ are $\MGL$-modules, we learn that when $E$ is one of these, then $\MGL \otimes E$ is a direct sum of pure suspensions of $E$. Again, \Cref{prop:sl-si-cond} finishes the job.
%
  Finally, for the case of $\kq_2$, we use the equivalence $\kq_2 \otimes C\eta \simeq \kgl_2$ and the fact that $\eta = 0$ in $\pi_{0,1,1} ^\R \MGL_2$ to deduce that $\MGL_2 \otimes \kgl_2 \simeq \MGL_2 \otimes \kq_2 \otimes C\eta \simeq \MGL_2 \otimes \kq_2 \oplus \Sigma^{1,1,1} \MGL_2 \otimes \kq_2$.
  It follows that $\MGL_2 \otimes \kq_2$ is a retract of $\MGL_2 \otimes \kgl_2$, so that $\MGL_2 \otimes \kq_2$ is slice simple.
\end{proof}
%
%
%\begin{cnstr}
%  The lax symmetric monoidal functor
%  \[\Gamma_* : \Sp_{C_2} \to \Sp_{C_2} ^{\Fil}\]
%  of SEC fators through a lax symmetric monoidal functor
%  \[\Gamma_* : \Sp_{C_2} \to \Mod(\Sp_{C_2}; \Gamma_* (\Ss)).\]
%  By REF, this may be identified with a lax symmetric monoidal funtor
%  \[\Gamma_* : \Sp_{C_2} \to \SHRat\]
%\end{cnstr}
%
%\begin{thm}
%  Suppose that $X \in \Sp_{C_2}$ admits a bounded below cell structure with cells of the form $\Ss^{n\rho}$. Then there is a natural equivalence
%  \[\nur(X) \otimes C\ta \simeq \pi_0 ^{\ch} \nur(X).\]
%\end{thm}
%
%The first main theorem of this section shows that $\nur$ is a decompletion of the functor
%\[\Gamma_* : \Sp_{C_2} \to \SHR^{\at}\]
%studied in SECTION.
%
%\begin{thm} \label{thm:Gamma-nu-comp}
%  For $X \in \Sp_{C_2}$, there is a natural equivalence
%  \[\Gamma_* (X) \simeq \nur (X) ^{\wedge} _{\MGL}.\]
%\end{thm}
%
%The second main theorem makes use of this new description of $\Gamma_*$ to prove that several $2$-complete $\R$-motivic spectra may be functorially recovered from their Betti realizations:
%
%\begin{thm}\label{thm:recover}
%  There are natural equivalences in $\SHRat$:
%  \[\MGL_2 \simeq \Gamma_* (\MU_{\R,2})\]
%  \[\HFt \simeq \Gamma_* (\uFt)\]
%  \[\HZt \simeq \Gamma_* (\uZt)\]
%  \[\kgl_2 \simeq \Gamma_* (\ku_{\R,2})\]
%  \[\kq_2 \simeq \Gamma_* (\ko_{C_2, 2}).\]
%\end{thm}
%
%\begin{rmk}
%  In light of \Cref{thm:recover}, it would be reasonable to define an $\R$-motivic spectrum of motivic modular forms as $\nur (\tmf_{C_2,2})$, assuming that one had a good equivariant spectrum of $C_2$-motivic modular forms $\tmf_C_2,2}$.
%  Since no such equivariant spectrum has yet been constructed, we leave this to future work. \todo{Mention why previous attempts don't work? Errors in the literature?}
%\end{rmk}
%


%% \subsection{The effective slice filtration} \label{sec:eff}
%% %The effective slice filtration is a useful computational tool in the study of motivic spectra. Moreover, Isaksen has suggested that $C_2$-spectra be studied by Betti realization of the effective slice filtration of $\R$-motivic spectra. Isaksen has called the associated spectral sequence the effective spectral sequence of the Betti realization.
%
%In this section, we give an explicit description of the effective slice filtration for Artin-Tate $\R$-motivic spectra in terms of the filtered model of \Cref{thm:filt-model}.

%Isaksen has suggested that $C_2$-spectra be studied via the Betti realization of the effective slice filtration of $\R$-motivic spectra. Isaksen has called the associated spectral sequence the $C_2$-effective spectral sequence. \todo{Is it a bad look that I attribute this to Dan even though the first place it appears is a paper of Hana's? What to do here?}
%As a consequence of our results, we observe that the $\ta$-Bockstein spectral sequence of an Artin-Tate $\R$-motivic spectrum contains equivalent data to the associated $C_2$-effective spectral sequence.

In this section, we relate the functor $i_*$ defined in the previous subsection to Voevodsky's effective slice filtration.
We begin by recalling the definition of the effective slice filtration \cite{VoeOpenProbs}.

\begin{dfn} \label{dfn:eff-slice}
  Let $\Sm / \R$ denote the category of smooth and separated $\R$-schemes of finite type.
  We let $\SHR^{\eff} _{i2, \geq n} \subset \SHR_{i2}$ denote the full subcategory generated under small colimits by the collection $\{\Ss_2^{p,q,q} \otimes X_+ \vert X \in \Sm / \R, p,q \in \Z, q \geq n\}$.
  We denote the right adjoint of the inclusion by $f_n : \SHR_{i2} \to \SHR^{\eff} _{i2, \geq n}$. There are natural transformations $f_{n+1} \to f_{n}$ and we let $s_n$ denote the cofiber of this map.

  Since $\SHR^{\eff} _{i2, \geq n} \subset \SHR_{i2}$ is a compactly generated stable subcategory, the functors $f_n$ and $s_n$ preserve all colimits. Moreover, the tensor product of an $n$-effective object with an $m$-effective object is $(m+n)$-effective, so that the effective slice tower functor $f_\bullet : \SHR_{i2} \to \SHR_{i2} ^{\Fil}$ is lax symmetric monoidal.
\end{dfn}

The main result of this section is the following:

\begin{prop} \label{thm:eff-fil}
  The lax symmetric monoidal functor $i_* : \SHRat \to \Sp_{C_2, i2}^{\Fil}$ of the previous subsection is equivalent to the lax symmetric monoidal functor $\b \circ f_\bullet : \SHRat \to \Sp_{C_2, i2}^{\Fil}$ taking $E \in \SHRat$ to the Betti realization of its effective slice tower
  \[\dots \to \b (f_{2} E) \to \b (f_1 E) \to \b (f_0 E) \to \b (f_{-1} E) \to \b (f_{-2} E) \to \dots.\]
\end{prop}
%
%\begin{rmk}
%  Let $X \in \SHRat$. Since $i_* (X \otimes C\ta) \simeq \Gr(i_* (X))$, the $\ta$-Bockstein spectral sequence is simply the usual spectral sequence computing the homotopy of the filtered object $i_* (X)$ from the homotopy of its associated graded.
%  On the other hand, Isaksen's $C_2$-effective spectral sequence starts from the homotopy of $\Gr(\b(f_* (X)))$ and attempts to compute the homotopy of $\b(X)$.
%
%  It follows from \Cref{thm:eff-fil} that, though the target of the $\ta$-Bockstein spectral sequence contains more information than the target of the $C_2$-effective spectral sequence, these spectral sequences ultimately contain the same information: their $\mathrm{E}_1$-pages and differentials may be read off from one another.
%\end{rmk}
%
%\begin{rmk}
%  In the case that $X \simeq \Gamma_* (Y)$, it follows that we obtain a purely topological description of the tower $\b (f_* (X))$. For example, we have $\kq_2 \simeq \Gamma_* (\ko_{C_2, 2})$ by \Cref{rmk:gamma-nu-same}; this case has been studied by Kong \cite{Kong}.
%\end{rmk}

As a first step, we rephrase \Cref{dfn:eff-slice} to be more natural in our trigraded context:

\begin{lem}
  The full subcategory $\SHR^{\eff} _{i2, \geq n} \subset \SHR_{i2}$ is generated under small colimits by the collection $\{\Ss_2 ^{p,q,w} \otimes X_+ \vert X \in \Sm / \R, p,q,w \in \Z, w \geq n\}$.
  As a consequence, the suspension functors provide equivalences,
  \[ \Sigma^{p,q,w} : \SHR^{\eff} _{i2, \geq n} \to \SHR^{\eff} _{i2, \geq n+w}. \]
  %% are equivalences of categories for each $p,q,w \in \Z$.
\end{lem}

\begin{proof}
  Since $\Sigma^{0,1,1} : \SHR^{\eff} _{i2, \geq n} \to \SHR^{\eff} _{i2, \geq n+1}$ is clearly an equivalence of categories, it suffices to prove the generation statement in the case that $n = 0$.
  Since $\SHR^{\eff} _{i2, \geq 0}$ is closed under tensor products, it suffices to show that $\Ss_2^{0,-1,0} \in \SHR^{\eff} _{i2, \geq 0}$.

  Now, $\SHR^{\eff} _{i2, \geq 0}$ is clearly closed under $\Sigma^{-1,0,0}$, so it suffices to show that $\Ss_2 ^{1,-1,0} \in \SHR^{\eff} _{i2, \geq 0}$.
  This follows from the existence of a cofiber sequence
  \[ \Ss^{0,0,0} \to \Spec \C \to \Ss^{1,-1,0}. \]
  The second statement is a clear consequence of the generation statement.
\end{proof}

We now define an alternative filtration on $\SHRat$ which is easier to analyze; following an argument of Heard \cite{SliceComp}, itself an adaptation of an argument of Pelaez \cite{Pelaez}, we will prove that this filtration is in fact equivalent to the effective slice filtration.

\begin{dfn}
  We let $\SHR^{\at, \gceff} _{i2, \geq n}$ denote the full subcategory generated under small colimits by the collection $\{\Ss_{2} ^{p,q,w} \vert p,q,w \in \Z, w \geq n\}$.
  We let
  \[ f_n ^{\att} : \SHRat \to \SHR^{\at, \gceff}_{i2, \geq n} \]
  denote the right adjoint of the inclusion. There are natural transformations $f_{n+1} ^{\att} \to f_{n} ^{\att}$, and we denote the cofiber by $s_n ^{\att}$.

  Since $\SHR^{\at, \gceff}_{i2, \geq n} \subset \SHRat$ is a compactly generated stable subcategory, the functors $f_n ^{\att}$ and $s_n ^{\att}$ preserve colimits.
\end{dfn}

It is clear that $\SHR^{\at,\gceff} _{i2, \geq n} \subset \SHRat \cap \SHR^{\eff} _{i2, \geq n}$.
%
%\begin{rmk}
%  Since the representation spheres $\Ss_2^{p+q\sigma}$ generate $\Sp_{C_2, i2}$ under colimits, $\SHR^{\at, \gceff} _{i2, \geq n}$ may be equivalently described as the full subcategory of $\SHRat$ generated under small colimits by the collection $\{ \Sigma^{0,0,w} c_{\C/\R} (\Sp_{C_2, i2}) \vert w \geq n\}$.
%\end{rmk}

\begin{lem} \label{lem:susp-slice}
  Given $E \in \SHR_{i2}$, there are natural equivalences $f_k \Sigma^{p,q,w} E \simeq \Sigma^{p,q,w} f_{k-w} E$ and $s_k \Sigma^{p,q,w} E \simeq \Sigma^{p,q,w} s_{k-w} E$.
  If $E \in \SHRat$, the analagous fact holds for $f_k ^{\att}$ and $s_k ^{\att}$.
\end{lem}

\begin{proof}
  This follows directly from the fact that $\Sigma^{p,q,w} : \SHR^{\eff} _{i2, \geq k} \to \SHR^{\eff}_{i2, \geq k +w}$ and $\Sigma^{p,q,w} : \SHR^{\at, \gceff} _{i2, \geq k} \to \SHR^{\at, \gceff}_{i2, \geq k +w}$ are equivalences of categories.
\end{proof}

\begin{lem} \label{prop:eff-slice-at-compare}
  Given  $E \in \SHRat$, there are natural equivalences $f_n E \simeq f_n ^{\att} E$ for all $n \in \Z$.
\end{lem}

\begin{proof}
  It is easy to see that the categories $\SHR^{\eff} _{i2, \geq n} \subset \SHR_{i2}$ and $\SHR^{\at, \gceff} _{i2, \geq n} \subset \SHRat$ define slice filtrations in the sense of \cite[Definition 2.1]{SliceComp}.
  The result will therefore follow from \cite[Theorem 2.20]{SliceComp} if we can verify three conditions. Let $\iota : \SHRat \hookrightarrow \SHR_{i2}$ denote the inclusion. Then we must show that the following hold for all $E \in \SHRat$:
  \begin{enumerate}
    \item The natural map $\varinjlim_{n} \iota(f_n ^{\att} E) \to \iota(\varinjlim_{n} f_n ^{\att} E))$ is an equivalence.
    \item $\iota(f_n ^{\att} E) \in \SHR^{\eff} _{i2, \geq n}$.
    \item $\Map_{\SHR_{i2} } (X, \iota(s_n ^{\att} E)) \simeq 0$ for all $X \in \SHR^{\eff} _{i2, \geq n+1}$.
  \end{enumerate}
  Condition (1) is clear from the fact that $\SHRat$ is closed under colimits in $\SHR_{i2}$, and condition (2) follows from the fact that $\SHR^{\at, \gceff} _{i2, \geq n} \subset \SHR^{\eff} _{i2, \geq n}$.

  To prove condition (3), we note that, since $s_n ^{\att}$ commutes with filtered colimits and $\SHR^{\eff} _{i2, \geq n+1}$ is compactly generated, it suffices to prove the statement for generators of $\SHR^{\at, \gceff} _{i2, \geq n}$, namely the trigraded spheres $\Ss_2 ^{p,q,w}$ where $w \geq n$.
  Since $s_n ^{\att} \Ss_2^{p,q,w} \simeq \Sigma^{p,q,w} s_{n-w} ^{\att} \Ss_2^{0,0,0}$ by \Cref{lem:susp-slice}, it suffices to show this for $\Ss_2^{0,0,0}$.

  This follows from the equivalence $s_{n} ^{\att} \Ss_2 ^{0,0,0} \simeq s_{n} \Ss_2^{0,0,0}$, which may be proved exactly as in \cite[Theorem 3.15]{SliceComp}.
\end{proof}

We are now free to use $f_n$ and $f^{\att} _n$ interchangeably. The following proposition gives us the needed control over $f^{\att} _n$:

\begin{prop} \label{prop:eff-condition}
  Let $E \in \SHRat$. Then $E \in \SHR^{\at, \gceff} _{i2, \geq n}$ if and only if
  \[ \pi^{\R} _{p,q,w} (C\ta \otimes E) = 0 \]
  for all $w < n$, i.e. if and only if $\ta : \pi^{\R} _{p,q,w+1} E \to \pi^{\R} _{p,q,w} E$ is an isomorphism for all $w < n$.
\end{prop}

\begin{proof}
  To show that $\pi^{\R} _{p,q,w} (C\ta \otimes E) = 0$ for all $w < n$ if $E \in \SHR^{\at, \gceff} _{i2,\geq n}$, it suffices to prove this when $E  = \Ss^{p,q,w}$ for $w \geq n$, which follows from \Cref{prop:weak-one-sided}.

  On the other hand, suppose that $E \in \SHRat$ satisfies $\pi^{\R} _{p,q,w} (C\ta \otimes E) = 0$ for all $w < n$.
  We will show that $f_n ^{\att} E \to E$ is an equivalence.
  First, we note that it is an equivalence on $\pi^{\R} _{p,q,w}$ for all $w \geq n$ by definition.
  By assumption, $\ta : \pi^{\R}_{p,q,w+1} E \to \pi^{\R} _{p,q,w} E$ is an isomorphism for all $w < n$.
  On the other hand, by the above $\ta : \pi^{\R}_{p,q,w+1} (f_n ^{\att} E) \to \pi^{\R} _{p,q,w}  (f_n ^{\att} E)$ is also an isomorphism for all $w < n$.
  This implies that $f_n ^{\att} E \to E$ in fact induces an isomorphism on all $\pi_{p,q,w} ^{\R}$, as desired.
\end{proof}
%
%\begin{cor} \label{cor:eff-slice-filt}
%  Let $E \in \SHRat$, and let
%  \[i_* (E) = \dots \to X_2 \to X_1 \to X_0 \to X_{-1} \to X_{-2} \to \dots.\]
%  On the other hand, let
%  \[i_* (f_{n} ^{\at} E) = \dots \to \overline{X}_2 \to \overline{X}_1 \to \overline{X}_0 \to \overline{X}_{-1} \to \overline{X}_{-2} \to \dots.\]
%  Then $\overline{X}_{w+1} \to \overline{X}_{w}$ is an equivalence for all $w < n$, and the map $i_* (f_n ^{\at} E) \to i_* (E)$ induces equivalences $X_w \xrightarrow{\sim} \overline{X}_w$ for $w \geq n$.
%\end{cor}
%
%\begin{proof}
%  This follows directly from \Cref{prop:eff-condition} and the definition of the functor $i_* : \SHRat \to \Fil(\Sp_{C_2,i2})$.
%\end{proof}

Finally, we are ready to prove \Cref{thm:eff-fil}.

\begin{proof}[Proof of \Cref{thm:eff-fil}]
  Given a bifiltered object $X_{\bullet, *}$, we let $\diag(X_{\bullet, *}) = X_{\bullet, \bullet}$ denote the filtered object obtained by restricting along the diagonal map $\Z^{\Fil} \hookrightarrow \Z^{\Fil} \times \Z^{\Fil}$.
  Then there is a span of lax symmetric monoidal functors:
  \begin{center}
    \begin{tikzcd}
      \diag \circ i_* \circ f_\bullet \ar[r] \ar[d] & i_* \\
      \b \circ f_\bullet, &
    \end{tikzcd}
  \end{center}
  where the horizontal map is induced by the natural transformation $f_\bullet \to Y$ and the vertical map is induced by the natural transformation $i_* \to Y \circ \b$. (Here as in \Cref{app:fil}, $Y$ is the functor taking an object to its constant filtered object.)

  Applied to $E \in \SHRat$, in filtration $n$ this span looks like
  \begin{center}
    \begin{tikzcd}
      \Map_{\SHRat} (\Ss_2 ^{0,0,n}, f_n E) \ar[r] \ar[d] & \Map_{\SHRat} (\Ss_2 ^{0,0,n}, E) \\
      \b (f_n E). &
    \end{tikzcd}
  \end{center}
  The horizontal map is an equivalence by $n$-effectivity of $\Ss_2 ^{0,0,n}$, and the vertical map is an equivalence by \Cref{thm:comparison-invert-ta} and \Cref{prop:eff-condition}.
\end{proof}
%
%\Cref{thm:eff-fil} now follows from combining \Cref{prop:eff-slice-at-compare} with \Cref{cor:eff-slice-filt}, in light of the fact that $\b (E) \simeq \varinjlim i_* (E)_n$.
%
%\begin{cor}
%  The functor $i_* : \SHR^{gc} _{i2} \to \Fil(\Sp_{C_2, i2})$ is equivalent to the functor $\b \circ f_* : \SHR^{gc} _{i2} \to \Fil(\Sp_{C_2, i2})$ taking $E \in \SHR^{gc} _{i2}$ to the tower
%  \[\dots \to \b f_{2} E \to \b f_1 E \to \b f_0 E \to \b f_{-1} E \to \b f_{-2} E \to \dots.\]
%\end{cor}
%
%\begin{rmk}
%  In the cases of interest where the effective spectral sequence has been studied
%\end{rmk}

